% RESEARCH_DRAFT_SHA256: df95208b018815107e727b060a7c8f4d676c899e1c1f63ecbc63dbc2d6d9eae1
% Layout and source-reference wrappers expanded for the web renderer.
\documentclass[11pt]{article}
\usepackage[margin=25mm]{geometry}
\usepackage{fontspec,unicode-math}
\setmainfont{Latin Modern Roman}
\setmathfont{Latin Modern Math}
\usepackage{amsmath,amssymb,mathtools,xurl,hyperref}
\hypersetup{colorlinks=true,urlcolor=blue,linkcolor=blue}
\setcounter{secnumdepth}{0}
\setlength{\parindent}{0pt}
\setlength{\parskip}{5pt}
\setlength{\emergencystretch}{3em}
\newcommand{\sref}[2]{\hyperlink{source:#1:#2}{[S#2]}}
\newcommand{\cataloguescope}[1]{\paragraph{Source record.}#1}
\newcommand{\floor}[1]{\left\lfloor#1\right\rfloor}
\begin{document}
\section*{Research-attempt provenance}
This standalone research attempt retains the complete original definition body below,
including its historical source and preparation notes. Those notes describe the
earlier catalogue preparation; the new attempt and its verification are identified
separately after that body. No original ID, catalogue status, or source statement is changed.

% UnsolvedMath ID 30006176; source catalogue code OWR-14299083-002
\section{Generalization Guarantees for Self-Supervised Representations}\label{30006176}
{\small UnsolvedMath ID 30006176 · Computer Science}
\subsection{Definitions and mathematical statement}

\paragraph{Shared notation.}Standard number systems are $\mathbb N,\mathbb Z,\mathbb Q,\mathbb R,\mathbb C$. Specialized notation and index ranges follow the cited source and the definitions below.


Self-supervised pretraining learns a representation from unlabeled observations and their specified augmentations. A downstream task fits a prediction rule using that representation; a generalization bound controls expected test loss or its deviation from empirical loss.
\paragraph{Record scope and qualifications.}
The catalogue title is retained as metadata. Named certificates, contracts, bounds or frameworks in that title are not asserted as results of the original source extract.
The augmentations, pretraining objective, downstream task family, sample scaling, and loss are not fixed in the supplied question.

Source-record limitation: The augmentation/pretraining objective, downstream task class, and generalization criterion are not specified.
\paragraph{Statement recorded in the supplied source.}
Establish statistical generalization guarantees across downstream prediction tasks for representations obtained by self-supervised pretraining on large unlabeled augmented datasets, beyond the presently understood kernel and infinite-width regimes. \href{https://doi.org/10.4171/owr/2025/4}{[S2]}
\paragraph{Source status.}
The supplied catalogue labels this source record partially solved; that metadata was checked against the supplied snapshot on 2026-10-11. The label is preserved for provenance and is not an independent verification of present open status or mathematical correctness. The catalogue literature assessment is dated 2026-08-22.
\subsection{Short English statement}
Prove downstream generalization guarantees for the stated pretrained representations beyond kernel and infinite-width limits.
\subsection{Sources}
\begin{itemize}
\item[S1] Ulam AI and the UnsolvedMath Contributors, UnsolvedMath: A Curated Collection of Open Mathematics Problems, record 30006176 (OWR-14299083-002). Catalogue curation and metadata: CC BY 4.0. The source extract is retained; definitions, English summaries and qualifications are editorial additions.
\url{https://huggingface.co/datasets/ulamai/UnsolvedMath}.
\item[S2] Overparametrization, Regularization, Identifiability and Uncertainty in Machine Learning (2025), Oberwolfach Reports, DOI 10.4171/owr/2025/4, pp. 175--208
\url{https://doi.org/10.4171/owr/2025/4}.
\item[S3] Oberwolfach Report 4/2025
\url{https://content.ems.press/assets/public/full-texts/serials/owr/22/1/14299083/online/10.4171-owr-2025-4.pdf}.
\item[S4] Esser--Fleissner--Ghoshdastidar, Non-parametric representation learning with kernels
\url{https://ojs.aaai.org/index.php/AAAI/article/view/29077}.
\item[S5] Anil--Esser--Ghoshdastidar, Wide contrastive models and kernel PCA
\url{https://arxiv.org/abs/2403.08673}.
\item Source-specific Oberwolfach Reports licences, including CC BY-SA 4.0 where applicable, are retained; see THIRD\_PARTY\_NOTICES.md.
\end{itemize}
\paragraph{Preparation and review.}
The original source and UnsolvedMath contributors retain attribution for the source extract. Codex (GPT-6; exact serving revision not exposed) assisted with definitions, English summaries, scope notes and formatting. No new human mathematical authorship or independently verified solution is asserted. Permanent source ID 30006176 and code OWR-14299083-002 checked against the complete supplied registry and current repository. Distinct upstream records, including repeated questions and broad or damaged extracts, are retained for complete coverage of the source CS category. A repeated record is not a new independent conjecture. Complete mathematical-equivalence and current-literature review has not been performed. Preparation and final review follow contribution prompts C and D, with source gaps disclosed. Structural and rendering checks do not establish a new theorem or open status.
\label{end:30006176}

\section{Literature and context check}
The consulted primary report, printed p. 197, identifies downstream generalization for self-supervised representations and the limitations of the kernel/initial-regime results discussed there. This is a broad research direction, not a fully quantified distribution-free theorem. The finite experiment below isolates assumptions any proposed universal guarantee would need.
\section{Research attempt: A downstream-label obstruction beyond representation injectivity}

Let $X$ be a finite input set with full-support distribution $P$, and let $\phi:X\to Z$ be a fixed representation. Every binary target $f:X\to\{0,1\}$ is a function of $\phi$ if and only if $\phi$ is injective. Sufficiency follows by defining the decoder at $\phi(x)$ to equal $f(x)$. For necessity, if $\phi(x)=\phi(y)$ for distinct points, choose $f(x)=0,f(y)=1$. Any decoder gives the same prediction distribution at these two points; its error contribution is at least $\min(P(x),P(y))$. Indeed, if its probability of predicting 1 is $q$, the contribution is $P(x)q+P(y)(1-q)$, minimized at an endpoint. On the uniform two-point distribution this lower bound is $1/2$.

Injectivity alone still supplies no label-efficiency theorem. Take $X$ uniform on $N\geq2$ points, let $n$ labelled samples be independent, and allow any unlimited unlabelled data and representation procedure independent of the labels. As an averaging device choose the $N$ target labels independently and fairly. At an independent test input not appearing in the labelled sample, its label is a fresh fair bit independent of the training data, representation, and learner randomness. Every prediction then errs with probability $1/2$. The probability that the test input was not seen is $(1-1/N)^n$. Therefore the average expected test error over the random target tables is at least
\[
\frac12(1-1/N)^n.
\]
Some deterministic target table attains at least this average risk. A uniform guarantee of expected risk at most $0<\eta<1/2$ for all target tables consequently requires
\[
n\geq \frac{\log(2\eta)}{\log(1-1/N)}.
\]
This reasoning remains valid if $\phi$ is the identity and all input probabilities are known exactly. It therefore rules out deriving universal label savings from unlabelled reconstruction or injectivity alone.

\paragraph{Augmentation route and its first gap.}
If an objective enforces $\phi(Ax)=\phi(x)$ for an augmentation $A$, then any downstream label not invariant on an $A$-orbit conflicts with exact decoding by the preceding collision argument. For example, on $X=\{0,1\}$ the flip augmentation forces a constant representation, while target $f(x)=x$ has risk at least $1/2$ under the uniform distribution. Thus a proposed transfer proof needs an assumption linking the augmentation invariances to the downstream task class. The unlabelled objective alone cannot establish that link.
\section{Verification and adversarial audit}
The target averaging argument ends with a deterministic worst target, not an assumption that actual labels are resampled at each example. Unlabelled data and random learner states remain independent of the chosen unseen labels. The denominator in the sample bound is negative, so the inequality reverses when logs are divided. The augmentation counterexample concerns deliberately incompatible tasks, not the empirical task distributions in the source.
\paragraph{Audit corrections.} Separated information preservation from statistical sample complexity and stated the independence of all unlabelled information from target labels.
\paragraph{External computation.} None; symbolic arguments and boundary checks only.
\section{Scope and remaining gaps}
No kernel-free generalization bound for a specified useful task family or training dynamics is proved. These are exact impossibility checks for an unrestricted interpretation, and the original research agenda remains open here.
\paragraph{Integration review.} Codex (GPT-6; exact serving revision not exposed) reviewed the written arguments and expanded layout and source-reference wrappers. Repository metadata and the production web renderer were checked using Python 3.13 and JavaScript with MathJax 3.2.2. This model-assisted review is not independent human proof verification; no PDF compilation or formal proof checking is asserted.

\section{Completion Estimate}
\noindent\textbf{COMPLETION: 0\%}
The calculations below the preserved source statement investigate a concrete route,
but establish no substantiated new progress toward the entire original target.
Known special cases and failed sufficient conditions are not credited as a solution.
\section{References}
The original source references and catalogue attribution are preserved in the
Sources section above. The following URLs were actually consulted for this attempt;
retention of an additional original reference does not claim fresh consultation of it.
\begin{itemize}
\item \url{https://ems.press/content/serial-article-files/51348} (consulted 2026-10-11). 
\end{itemize}
\section{Model and contributor disclosure}
Codex GPT-6 completed written-proof review of this attempt at its declared scope; this is model-assisted review, not independent human proof checking.
\paragraph{Material model collaboration.} Root Codex GPT-6 supplied the unrestricted downstream-label sampling lower bound and finite representation-collision route. The gpt-6-astra author inspected and reworked these arguments; this is model-assisted mathematical work, not independent human verification.
OpenAI gpt-6-astra, requested reasoning configuration ultra, generated this research
attempt on 2026-10-11. The exact serving revision is not exposed. Codex GPT-6
performed orchestration and repository-format review; its exact serving revision
is also not exposed. The mathematical route, boundary checks and final adversarial
self-audit followed contribution prompts A then B in
\texttt{docs/contribution\_prompts.md}, after reading \texttt{CONTRIBUTING.md}.
Python 3.14.2 assembled the standalone documents; any mathematical computation is
specified in Verification. No theorem prover was used. No human mathematical
authorship, substantial human proof correction, or independent human proof checking
is asserted. Model-assisted review, source retrieval, and structural checks do not
establish correctness or novelty. Full private reasoning transcripts are unavailable
as reproducibility artifacts; the complete arguments, inputs and outcomes actually
used are stated here or in the companion batch audit.
\section{Final status}
Unresolved as a new research contribution. This is a concrete route investigation,
not a claim of new complete resolution or a current-status certification. A known
answer reproduced for context does not become a new contribution. The catalogue's
source-status label is preserved; this attempt status does not assert that a
source-labelled solved question is open.
\end{document}
